\section{A sharp spanning interface}
\label{sec:sharp-spanning}
\subsection{The exact source hypotheses}

Let \(S,T\) be finite-dimensional complex Hilbert spaces, with
\(s=\dim S>0\). Let a finite group \(G\) act unitarily on \(T\), by
\(\rho\), and let \(P\) be the orthogonal projection onto a
\(G\)-stable subspace \(W\), of dimension \(d>0\). Thus
\(P\rho(g)=\rho(g)P\). Fix a linear map \(B:S\to T\) and write
\(C_g=P\rho(g)B:S\to W\).

The following are the two quantitative hypotheses underlying
\texttt{isSpread\_of\_second\_moment}, in
\src{New/Spanning/SecondMoment/FewTranslates.lean}{\texttt{FewTranslates.lean}}:
for every subspace \(U\le S\),
\begin{equation}\label{eq:near}
 \|PB P_U\|_\HS^2\ge \dim U-u s,
\end{equation}
and
\begin{equation}\label{eq:moment}
 \frac d{|G|}\sum_{g\in G}\|B^*\rho(g)B\|_\HS^2
 \le (1+v)s^2.
\end{equation}
Here \(P_U\) denotes orthogonal projection onto \(U\), and
\(0\le u<1\) and \(v\ge0\). The source's \texttt{IsNearIsometry} is
\eqref{eq:near}; it is a statement on every source subspace, not merely
one lower bound on the total Hilbert--Schmidt norm.

\begin{lemma}[The square-root estimate]\label{lem:sqrt}
For every subspace \(F\le W\), of relative dimension
\(\alpha=\dim F/d\), some \(g\in G\) satisfies
\begin{equation}\label{eq:sqrt}
 \dim C_g^{-1}(F)
 \le s\bigl[u+\sqrt{(1+v)\alpha}\bigr].
\end{equation}
In particular, \(s\le d/(1-u)\).
\end{lemma}

\begin{proof}
Put
\(M=|G|^{-1}\sum_g\rho(g)BB^*\rho(g)^*\).
Unitarity and a change of variables in the double sum give
\[
 \tr(M^2)=\frac1{|G|}\sum_g\|B^*\rho(g)B\|_\HS^2.
\]
The Hilbert--Schmidt Cauchy--Schwarz inequality therefore gives
\[
 \frac1{|G|}\sum_g\|P_F\rho(g)B\|_\HS^2
 =\tr(P_FM)
 \le\sqrt{\dim F\,\tr(M^2)}
 \le s\sqrt{\alpha(1+v)}.
\]
Choose \(g\) with its summand at most this average and put
\(U=C_g^{-1}(F)\). On \(U\), the maps \(C_g\) and
\(P_F\rho(g)B\) agree. Because \(C_g=\rho(g)PB\), and
right-composition by an orthogonal projection does not increase the
Hilbert--Schmidt norm, \eqref{eq:near} yields
\[
 \dim U-u s
 \le\|PB P_U\|_\HS^2
 =\|C_gP_U\|_\HS^2
 \le\|P_F\rho(g)B\|_\HS^2.
\]
This proves \eqref{eq:sqrt}. For the last assertion, apply
\eqref{eq:near} to \(\ker(PB)\). It gives
\(\dim\ker(PB)\le u s\), hence
\((1-u)s\le\operatorname{rank}(PB)\le d\).
\end{proof}

The source first replaces both errors by a common \(\varepsilon\), then applies
\(\sqrt{\alpha(1+\varepsilon)}\le(\alpha+1+\varepsilon)/2\)
before proceeding to the greedy argument. Keeping the square root
retains useful information near the limiting dimension.

\subsection{Sharp total span and a finite selection bound}

Set
\begin{equation}\label{eq:parameters}
\begin{gathered}
 a=1-u,\qquad z=1+v,\qquad
 \alpha_*=\frac{a^2}{z},\\[3pt]
 \delta_*=1-\alpha_*=
 \frac{2u+v-u^2}{1+v},
 \qquad h=\frac{s z}{2ad},\qquad q=\max\{0,1-h\}.
\end{gathered}
\end{equation}
Thus \(a,z,\alpha_*,h>0\), \(0\le\delta_*<1\), and
\(0\le q<1\).

\begin{theorem}[Greedy selection with the sharp limiting loss]\label{thm:greedy}
For every integer \(t\ge0\), there is a set \(A\subseteq G\),
with \(|A|\le t\), such that
\begin{equation}\label{eq:greedy}
 \codim_W\sum_{g\in A}\im C_g
 \le d\bigl[\delta_*+\alpha_*q^t\bigr].
\end{equation}
Consequently, for \(0<\tau<\alpha_*\), a loss of at most
\((\delta_*+\tau)d\) is achieved with at most
\begin{equation}\label{eq:count}
 \left\lceil\frac{2ad}{sz}\log\frac{\alpha_*}{\tau}\right\rceil
\end{equation}
translates. The span of all the images has codimension at most
\(\lfloor\delta_*d\rfloor\).
\end{theorem}

\begin{proof}
Let the current span be \(F\), with codimension \(c\). Rank-nullity
for \(S\to W/F\) gives
\[
 \dim(F+\im C_g)-\dim F=s-\dim C_g^{-1}(F).
\]
By Lemma~\ref{lem:sqrt}, some \(g\) therefore adds at least
\begin{equation}\label{eq:gain}
 s\left[a-\sqrt{z(1-c/d)}\right]
\end{equation}
dimensions. Suppose \(c>\delta_*d\), and put
\(e=c-\delta_*d>0\). The bracket in \eqref{eq:gain} is
\[
 a-\sqrt{a^2-ze/d}
 =\frac{ze/d}{a+\sqrt{a^2-ze/d}}
 \ge\frac{ze}{2ad}.
\]
The square root is real because \(c\le d\). Thus, after the
selected step, the new excess codimension is at most \((1-h)e\),
and hence at most \(qe\). If \(c\le\delta_*d\), keep the
existing set of indices. Induction from the empty span,
\(c_0=d=\delta_*d+\alpha_*d\), proves \eqref{eq:greedy}.
Taking a set of indices accommodates repetitions automatically.

Since \(q\le e^{-h}\), the number in \eqref{eq:count} suffices.
This also covers \(q=0\), as the displayed ceiling is at least one.
Finally, take \(F=\sum_{g\in G}\im C_g\) in
\eqref{eq:sqrt}. Every preimage is all of \(S\), so
\(1\le u+\sqrt{z\dim F/d}\), whence
\(\dim F\ge\alpha_*d\). Integrality gives the stated floor.
\end{proof}

\begin{proposition}[Sharpness under the moment hypotheses]\label{prop:sharp}
For every rational pair \(0\le u<1\), \(v\ge0\), examples
satisfying \eqref{eq:near} and \eqref{eq:moment} have total span of
dimension exactly \(\alpha_*d\). Thus \(\delta_*\) cannot be
decreased as a universal bound under these hypotheses.
\end{proposition}

\begin{proof}
Since \(\alpha_*\) is positive and rational, choose positive integers
\(d,r\) with \(r/d=\alpha_*\). Take \(S=\C\),
\(W=T=\C^d\), and \(P=I\). Let the cyclic group of order \(r\)
rotate the first \(r\) coordinate vectors, fixing the other
\(d-r\) vectors. Put \(B(1)=\sqrt{1-u}\,e_1\).
The only subspaces of \(S\) are zero and \(S\), so
\eqref{eq:near} holds, with equality for \(S\). Moreover
\[
 \frac1{|G|}\sum_g\|B^*\rho(g)B\|_\HS^2
 =\frac{(1-u)^2}{r},\qquad
 \frac{d(1-u)^2}{r}=1+v.
\]
Thus \eqref{eq:moment} is also an equality. The orbit images span
exactly the first \(r\) coordinate directions.
The example also has \(\|B^*B-I\|_\HS^2=u^2s\).
\end{proof}

\begin{corollary}[Equal errors]
When \(u=v=\varepsilon\), the sharp limiting deficit is
\[
 \delta_*=\frac{3\varepsilon-\varepsilon^2}{1+\varepsilon}.
\]
For \(\varepsilon=p/v_0\), with integers \(0\le p<v_0\),
one explicit equality example has
\(d=v_0(v_0+p)\) and \(r=(v_0-p)^2\).
\end{corollary}

The total-rank argument is an instance of the classical
second-moment, or effective-rank, inequality \cite{welch,mixon}.
Greedy rank accumulation is likewise a classical technique; see, for example,
the submodular-cover analysis in \cite{wolsey}. The contribution here
is the exact sharp threshold extracted from the source hypotheses,
the explicit selection schedule, and its application with all of
the source's rank-loss and translate-count terms preserved.

\subsection{Restriction of the source and its cost}

\begin{lemma}[Restriction loss]\label{lem:restriction}
Let \(S_0\le S\) have codimension \(r_0\). The same selected
set in Theorem~\ref{thm:greedy} satisfies
\begin{equation}\label{eq:restricted}
 \codim_W\sum_{g\in A}C_g(S_0)
 \le d\bigl[\delta_*+\alpha_*q^t\bigr]+t r_0.
\end{equation}
\end{lemma}

\begin{proof}
The direct sum of \(|A|\) copies of \(S/S_0\) surjects onto
\((\sum_{g\in A}C_g(S))/(\sum_{g\in A}C_g(S_0))\).
Thus the difference of the two span dimensions is at most
\(|A|r_0\le t r_0\). Apply \eqref{eq:greedy}.
\end{proof}

For an arbitrary application, \eqref{eq:restricted} explicitly
balances the benefit of selecting more translates against the
cost of restricting the source. If \(r_0/s=\theta>0\) and
\(2\theta<a\), taking
\[
 t=\left\lceil h^{-1}\log\frac{a}{2\theta}\right\rceil
\]
gives the further bound
\[
 \frac1d\codim_W\sum_{g\in A}C_g(S_0)
 \le\delta_*+
 \frac{2a\theta}{z}\left(1+\log\frac{a}{2\theta}\right)
 +\frac{\theta s}{d}
 \le\delta_*+
 \frac{2a\theta}{z}\left(1+\log\frac{a}{2\theta}\right)
 +\frac{\theta}{a}.
\]
Indeed, \(d\alpha_*h=as/2\), and rounding up increases the
linear loss by at most \(r_0\). This optimized choice is optional;
the next theorem improves the original source's count budget.

\section{Improving the source theorem at its original constants}
\label{sec:source-improvement}

The parameter called \(b\) here is the large weight in the source,
not the moment factor \(z=1+v\). The raw estimates in
\src{New/Spanning/SecondMoment/SpanningTheorem.lean}{\texttt{SpanningTheorem.lean}}
are
\[
 u=\frac5{b^4},\qquad v=\frac9{b^4}.
\]
They occur in the following declarations, respectively:
\begin{itemize}
\item \code{isNearIsometry_e}\textsubscript{0}\code{_comp_scaledPoincare};
\item \code{targetDim_mul_sum_hilbertSchmidtNorm_sq_le}.
\end{itemize}
The proof of \texttt{isSpread\_goodPoincare} later replaces both
by \(\varepsilon=(8b^2)^{-1}\); that enlargement is unnecessary
for the argument here. Its source-restriction argument in
\src{New/Reduction/ShortAndEnough.lean}{\texttt{ShortAndEnough.lean}}
has \(r_0/s\le b^{-8}\). The maps are identified through
\texttt{sourceFormCoords}; the nonzero scalar normalizations in
\texttt{goodPoincare} and \texttt{normalizedPoincare} do not change
these image subspaces.

\begin{theorem}[Fourth-order loss and fewer translates]\label{thm:source}
Assume \eqref{eq:near} and \eqref{eq:moment} with
\(u=5b^{-4}\), \(v=9b^{-4}\), and \(b\ge6\).
Let \(S_0\le S\) satisfy \(r_0=s-\dim S_0\le b^{-8}s\).
Then some set \(A\subseteq G\) satisfies
\begin{align}
 |A|\dim S_0&\le 4d\log(1+2b^2),\label{eq:sourcecount}\\
 \dim\sum_{g\in A}C_g(S_0)&\ge
 \left(1-\frac{20}{b^4}\right)d.\label{eq:sourcerank}
\end{align}
In fact \(|A|s\le4d\log(1+2b^2)\).
\end{theorem}

\begin{proof}
Write \(L=\log(1+2b^2)\), and choose
\[
 t=\left\lfloor\frac{4dL}{s}\right\rfloor.
\]
Lemma~\ref{lem:sqrt} gives \(s/d\le1/a\). Since
\(u\le5/256\) and \(v\le9/256\), we have
\(a\ge251/256\), \(z\le265/256\), and
\[
 0<h=\frac{s z}{2ad}\le\frac{z}{2a^2}
 \le\frac{33920}{63001}<1.
\]
Also \(4dL/s\ge4aL>1\), so \(t\ge1\).
In particular \(q=1-h\), and
\begin{align*}
 q^t&\le e^{-ht}
 \le \exp\left(\frac{s z}{2ad}\right)
       (1+2b^2)^{-2z/a}\\
 &\le \exp\left(\frac{z}{2a^2}\right)
       (1+2b^2)^{-2z/a}.
\end{align*}
The second inequality uses \(t\ge4dL/s-1\); this retains the
entire cost of the floor. Applying Lemma~\ref{lem:restriction}
and dividing by \(d\), the missing share is at most
\begin{equation}\label{eq:exactsource}
 \delta_*+
 \alpha_*\exp\left(\frac{z}{2a^2}\right)
       (1+2b^2)^{-2z/a}
 +4b^{-8}\log(1+2b^2).
\end{equation}
Bound the three terms separately. First,
\[
 \delta_*=\frac{2u+v-u^2}{1+v}\le2u+v=\frac{19}{b^4}.
\]
Second, \(\alpha_*\le1\), \(z/(2a^2)<1\), \(e<3\),
and \(z/a\ge1\), so the middle term is at most
\[
 3(1+2b^2)^{-2}\le\frac3{4b^4}.
\]
Third, \(\log(1+2b^2)\le2b^2\), so the final term is at most
\(8b^{-6}\). For \(b\ge6\),
\[
 \frac{19}{b^4}+\frac3{4b^4}+\frac8{b^6}
 \le\frac1{b^4}\left(19+\frac34+\frac8{36}\right)
 =\frac{719}{36b^4}<\frac{20}{b^4}.
\]
This proves \eqref{eq:sourcerank}. Finally,
\(|A|\dim S_0\le |A|s\le ts\le4dL\), proving
\eqref{eq:sourcecount}.
\end{proof}

\begin{corollary}[Original budget, smaller minimum weight]\label{cor:five}
For \(b\ge4\), the conclusion \eqref{eq:sourcerank} holds with
\(|A|s\le5d\log(1+2b^2)\).
\end{corollary}
\begin{proof}
Take \(t=\lfloor5dL/s\rfloor\). The middle term in
\eqref{eq:exactsource} becomes at most
\(3(1+2b^2)^{-5/2}\le b^{-5}\), and the restriction term is
at most \(10b^{-6}\). The total is at most
\(b^{-4}(19+1/4+10/16)=159/(8b^4)<20/b^4\).
\end{proof}

\begin{proposition}[The fourth-order scale and leading coefficient are optimal]
For every integer \(b\ge2\), the raw error bounds
\(u=5b^{-4}\), \(v=9b^{-4}\), with \(S_0=S\), admit an
example whose missing share even after all translates is
\[
 \frac{19b^4-25}{b^4(b^4+9)}
 =\frac{19}{b^4}-\frac{196}{b^4(b^4+9)}.
\]
Consequently no \(o(b^{-4})\) universal deficit bound, or
universal leading coefficient below \(19\), follows from these
raw moment and near-isometry inequalities alone.
\end{proposition}
\begin{proof}
Use Proposition~\ref{prop:sharp} with
\(d=b^4(b^4+9)\), \(r=(b^4-5)^2\), and \(S=\C\).
Then \(d-r=19b^4-25\) and \(r/d=(1-u)^2/(1+v)\).
The unrestricted source meets the ramification-loss bound with
zero loss. The stated asymptotic follows by division.
\end{proof}

\begin{remark}[Precise status and significance]
The existing \texttt{enough\_forms} in the pinned source has
\((1-4/b^2)d\) in place of \eqref{eq:sourcerank}, with the count constant five in place of four in \eqref{eq:sourcecount}. The source imposes much
larger lower bounds on \(b\) than \(b\ge6\), so this is an
improvement throughout its admissible parameter range. Theorem~\ref{thm:source}
uses the original moment and ramification estimates; it requires no
new arithmetic input. All three errors, the floor, and the restriction
of the source have been retained. The count constant is reduced from five to four; Corollary~\ref{cor:five} retains five when only \(b\ge4\) is assumed.

This is an ordinary mathematical proof of the strengthened statement.
The upstream code has not been modified, and the new statement has not
been executed through Lean. It improves the tolerance for completing
an almost-full span to a basis: with \(O(b^{-4})\) missing dimensions,
the per-missing-vector completion cost can be much larger for a fixed
total budget than with \(O(b^{-2})\) missing dimensions. It does not
by itself change the \(x^{-1/2+o(1)}\) prime-sum exponent or establish
another arithmetic RH.
\end{remark}

